\documentclass[a4paper]{article}
%\usepackage[swedish]{babel}
\usepackage[utf8]{inputenc}
\usepackage{a4wide}
\usepackage{appendix}
%vid problem, testa någon av följande:
%\usepackage[cp1252]{inputenc} under Windows
%\usepackage[applemac]{inputenc} under MacOS 9 eller tidigare
%\usepackage[latin1]{inputenc} under MacOS X

\usepackage[T1]{fontenc}
\usepackage[pdftex]{graphicx}

%\LaTeX\ är kul. Den senaste stora versionen heter \LaTeX2e och släpptes 1994~\cite{lamport94}.
%\begin{equation}
%a^2+b^2=c^2
%\label{eq:Pythagoras}
%\end{equation}
%För Pythagoras ekvation, se Ekvation~\ref{eq:Pythagoras}.

\title{A* Pathfinding Acceleration with use of Auto-Generated Waypoints for Grid Traversal}
\author{Fredrik Olsson, Magnus Nyqvst}
\date{\today} 

\begin{document}
\pagenumbering{gobble}
\maketitle
\newpage
\thispagestyle{empty}
\paragraph{Abstract--}
Sammanfattar rapporten
Varför är vår rapport värd att läsa?
Syfte, metod
Viktiga resultat och slutsatser Nyckelord
“Tänk på att detta skall kunna läsas fristående”

\tableofcontents
\listoffigures
\newpage
\pagenumbering{arabic}
\twocolumn
\section{Introduction}
Pathfinding is a fundamental part of games \cite{dynaPF15}\cite{roboGame15} and it is often supplemented by a waypoint graph to make traversal of a given region easier \cite{dynaPF15}. Every node in a waypoint graph is called a waypoint and they represent key locations in the region \cite{dynaPF15}. Each waypoint has edges towards other waypoints to where an object can travel through without risk of colliding with the surroundings \cite{dynaPF15}.\newline
In this paper, we propose a method to reduce execution time of the A* pathfinding algorithm. We improve our previously implemented A* algorithm with automatically generated and connected waypoints in a two-dimensional grid coordinate system. The waypoint generation is done in two steps. First, we generate a waypoint for each corner of an obstacle. Second, we check connections for every waypoint by sending a ray towards all other waypoints in the region. The waypoints are connected if the rays path is unblocked. Our waypoint generation method is heavily influenced by the one suggested in the work of Weiping et al. \cite{dynaPF15}.\newline
Executing pathfinding in dynamic environments is more challenging than in static environments \cite{dynaPF15}, and this study is therefore limited to completely static environments. The difference between the two terms path and shortest path is significant \cite{heuristicGame15}. We conducted studies of several pathfinding combinations, with and without waypoints, but we decided not to measure the time consumption of the shortest path with only A* pathfinding.



This is introduction lol \newline
Syfte \newline
Frågeställning \newline
Hypoteser \newline
Avgränsningar

\section{Background}
Stort spel \newline
Bakgrundsfakta (Definitioner som du använder dig av senare) \newline
Saker som läsaren behöver veta för att förstå din rapport \newline
Vad har gjorts tidigare?

\section{Related work}
Informationssökning \newline
Urval av litteratur - skriv inte om allt utan det viktigaste för din rapport

\section{Method}
Hur vi autogenerar (mycket referens till research articles) \newline
Vad har du gjort? \newline
Vilka metoder har du använt? \newline
Vad har du kommit fram till? \newline
Var noggrann och utförlig så att det går följa vad du gjort? \newline
Motivera om du antar något \newline
Diskutera begränsningar

\begin{figure}[h]
\centering
\includegraphics[scale=0.9, angle=20]{bth_logo.jpg}
\includegraphics[scale=0.8, angle=40]{bth_logo.jpg}
\includegraphics[scale=0.7, angle=60]{bth_logo.jpg}
\caption{Roterade BTH-logotyper}
\label{fig:Logo}
\end{figure}

\section{Result}
Bevisa att auto waypoints äger fett

Bilder kan roteras med hjälp av en dator, se exempel i Figur~\ref{fig:Logo}.

\section{Discussion}
Discuss the results

\begin{thebibliography}{9}

\bibitem{dynaPF15}
  Weiping Zhu, Daoyuan Jia, Hongyu Wan, Tuo Yang, Cheng Hu, Kechen Qin, Xiaohui Cui,
  \textit{Waypoint Graph Based Fast Pathfinding in Dynamic Environment},
  International Journal of Distributed Sensor Networks, vol. 2015, Article ID 238727, 12 pages,
  August 2015. https://doi.org/10.1155/2015/238727.

\bibitem{roboGame15}
  Zeyad Abd Algfoor, Mohd Shahrizal Sunar, Hoshang Kolivand,
  \textit{A Comprehensive Study on Pathfinding Techniques for Robotics and Video Games},
  International Journal of Computer Games Technology, vol. 2015, Article ID 736138, 11 pages,
  2015. https://doi.org/10.1155/2015/736138.

\bibitem{heuristicGame15}
  Geethu Elizebeth Mathew,
  \textit{Direction Based Heuristic For Pathfinding In Video Games},
  Procedia Computer Science, vol. 47, pp. 262-271,
  2015.

\bibitem{lamport94}
  Leslie Lamport,
  \textit{\LaTeX: a document preparation system},
  Addison Wesley, Massachusetts,
  2nd edition,
  1994.

\end{thebibliography}
\appendix
\cleardoublepage
\section{Introduction}
\section{Background}
\section{Related work}
\section{Method}
\section{Result}
\section{Discussion}
\end{document}
